\documentclass[a4paper,11pt]{article}

\usepackage[utf8]{inputenc}
\usepackage{amsmath,amssymb}
\usepackage[a4paper,margin=2.5cm]{geometry}
\usepackage{enumitem}

\title{Trigonometric functions}
\author{}
\date{}

\begin{document}

\maketitle

\section*{Introduction: Prior knowledge of trigonometry}

Draw a right-angled triangle on the board and mark one angle $\alpha$. Students in groups, they discuss questions

Ask the students:

\begin{itemize}
    \item What do you already know about sine and cosine?
    \item Where have you seen these functions before?
    \item What is special about the triangle we need?
    \item Which side is the hypotenuse?
    \item Which sides are the opposite and adjacent sides with respect to
    the angle $\alpha$?
\end{itemize}
\subsection*{1. Definitions}
Write definitions on the board :

\begin{equation}
    \sin(\alpha)
    =
    \frac{\text{opposite side}}{\text{hypotenuse}}
\end{equation}

\begin{equation}
    \cos(\alpha)
    =
    \frac{\text{adjacent side}}{\text{hypotenuse}}
\end{equation}

.


\subsection*{2. Similar Triangles}

Hand out \textbf{Exercise 1} or draw different triangles, some of them are similar. 

Let them compare triangles that have the same angle but different sizes.

\textbf{Discussion:}

\begin{itemize}
    \item What is different between the triangles?
    \item What stays the same?
    \item What happens to the side lengths when we make the triangle bigger?
    \item What happens to the ratios of the side lengths?
    \item What do you notice about $\sin(\alpha)$ and $\cos(\alpha)$?
\end{itemize}

The students should discover that similar triangles have different side
lengths, but the ratios of the corresponding sides remain the same.

Therefore,

\begin{center}
    \fbox{\parbox{0.8\textwidth}{
    \centering
    For a fixed angle $\alpha$, the values of
    $\sin(\alpha)$ and $\cos(\alpha)$ are always the same.

    They depend on the \textbf{angle}, not on the size of the triangle.
    }}
\end{center}

Ask:

\begin{quote}
    ``So if I tell you only the angle $\alpha$, do you need to know
    which particular triangle I am talking about in order to know
    $\sin(\alpha)$?''
\end{quote}

This motivates thinking of sine and cosine as \textbf{functions of an angle}:

\begin{equation}
    \alpha
    \longmapsto
    \sin(\alpha),
    \qquad
    \alpha
    \longmapsto
    \cos(\alpha).
\end{equation}


\subsubsection*{3. A Useful Relationship}

Continue with \textbf{Exercise 2}.

The students use Pythagoras' theorem to determine the missing side and
calculate $\sin(\alpha)$ and $\cos(\alpha)$.

Then ask them to calculate

\begin{equation}
    \sin^2(\alpha)+\cos^2(\alpha).
\end{equation}

Ask:

\begin{quote}
    ``Do you think the result is just a coincidence?''
\end{quote}

Let the students try to explain their observation using the triangle.

If the opposite side is $a$, the adjacent side is $b$ and the hypotenuse
is $c$, then

\begin{equation}
    \sin^2(\alpha)+\cos^2(\alpha)
    =
    \frac{a^2}{c^2}+\frac{b^2}{c^2}
    =
    \frac{a^2+b^2}{c^2}.
\end{equation}

Using Pythagoras' theorem,

\begin{equation}
    a^2+b^2=c^2,
\end{equation}

we obtain

\begin{equation}
    \boxed{\sin^2(\alpha)+\cos^2(\alpha)=1}.
\end{equation}

\textbf{this correlation is not coincidence but holds for all angles }

\subsubsection*{4. Short Application}

Use \textbf{Exercise 3} as a short application.

A ladder of length $5\,\mathrm{m}$ leans against a wall at an angle of
$60^\circ$.

Ask:

\begin{itemize}
    \item Which side of the triangle do we know?
    \item Which side do we want to determine?
    \item Should we use sine or cosine?
    \item Can you write down an equation for the height $h$?
\end{itemize}

The students should obtain

\begin{equation}
    \sin(60^\circ)=\frac{h}{5\,\mathrm{m}}
\end{equation}

and therefore

\begin{equation}
    h=5\,\mathrm{m}\cdot\sin(60^\circ).
\end{equation}


\subsubsection*{5. Transition to the Unit Circle}

Finish this part with a question that exposes a limitation of the triangle
definition.

Write on the board:

\begin{equation}
    \sin(30^\circ)
\end{equation}

and ask:

\begin{quote}
    ``Can we understand this using a right-angled triangle?''
\end{quote}

Then write:

\begin{equation}
    \sin(120^\circ).
\end{equation}

Ask:

\begin{quote}
    ``What about this one? Can a right-angled triangle have an angle of
    $120^\circ$?''
\end{quote}

The answer is no.

Nevertheless, a calculator gives a perfectly well-defined value for
$\sin(120^\circ)$ and even solution for angles $\alpha > 360 ^\circ$.

\textbf{Guiding question:}

\begin{center}
    \fbox{\parbox{0.8\textwidth}{
    \centering
    How can we extend our idea of sine and cosine so that they work
    for \textbf{any angle}, not only angles inside a right-angled triangle? 
    }}
\end{center}

If no one has an answer, name \textbf{unit circle}.

\subection*{ The Unit circle}

\begin{quote}
    ``Instead of putting our angle inside a triangle, we can now think of the angle as a rotation in a coordinate system.''
\end{quote}

\begin{itemize}
\item Draw circle with radius 1, a straight line through the origin and the circle in the first quadrant 
\item endpoint coordinates P(x,y) 
\item for the first quadrant: draw corresponding right-angled triangle 
\end{itemize}

\textbf{Questions:}

\begin{itemize}
    \item What is the length of the hypotenuse?
    \item What is the adjacent side?
    \item What is the opposite side?
    \item calculate sine and cosine of the angle 
\end{itemize}

Since radius of the circle is $1$ 

\begin{equation}
    \cos(\alpha)
    =
    \frac{x}{1}
    =
    x
\end{equation}

and

\begin{equation}
    \sin(\alpha)
    =
    \frac{y}{1}
    =
    y.
\end{equation}

Therefore,

\begin{equation}
    \boxed{P=(\cos\alpha,\sin\alpha)}.
\end{equation}

\begin{itemize}
    \item not a completely new definition
    \item for angles $ < 90^\circ $ definition equals our former definition with a right triangle 
\end{itemize}


\subsubsection*{ Going Beyond $90^\circ$}

Now rotate the radius further around the circle.

Consider, for example,

\begin{equation}
    \alpha=120^\circ.
\end{equation}

Ask:

\begin{itemize}
    \item Where is the point on the circle?
    \item Is its $x$-coordinate positive or negative?
    \item Is its $y$-coordinate positive or negative?
    \item What does this tell us about $\cos(120^\circ)$?
    \item What does this tell us about $\sin(120^\circ)$?
\end{itemize}

For the second quadrant we notice:

\begin{equation}
    \cos(\alpha)<0,
    \qquad
    \sin(\alpha)>0.
\end{equation}

Continue briefly through the other quadrants and let the students
predict the signs of sine and cosine.

\begin{center}
\begin{tabular}{c|c|c}
    \textbf{Quadrant} & $\boldsymbol{\cos(\alpha)}$
                      & $\boldsymbol{\sin(\alpha)}$ \\ \hline
    I   & $+$ & $+$ \\
    II  & $-$ & $+$ \\
    III & $-$ & $-$ \\
    IV  & $+$ & $-$
\end{tabular}
\end{center}

\subsubsection*{ Degrees and Radians}

There is another way of measuring angles called radians
\paragraph{An intuitive way to think about radians}

Imagine a circle with radius $1$. Start at the point $(1,0)$ and move
along the edge of the circle.

Instead of measuring how far you have turned in degrees, we can measure
the angle by asking:

\begin{quote}
    ``How far have I travelled along the circle until I reach the point associated with the angle?''
\end{quote}

For a circle with radius $1$, this distance is exactly the angle measured
in \textbf{radians}.

The circumference of the unit circle is

\begin{equation}
    C = 2\pi.
\end{equation}

Therefore, after travelling once around the entire circle, we have
travelled a distance of $2\pi$. This corresponds to a full rotation:

\begin{equation}
    360^\circ = 2\pi.
\end{equation}

Can you imagine what the following angles are in radians:
\begin{itemize}
    \item   $180^\circ = \pi$
    \item  $90^\circ = \frac{\pi}{2}$
    \item $270^\circ = \frac{3\pi}{2},$
    \item $360^\circ = 2\pi.$
\end{itemize}
    evtl. \textbf{Exercise 4}

\subsection*{5. Exploring the Unit Circle}

evtl.  \textbf{Exercise 5}.

Let the students label the four points on the unit circle and complete
the table for

\begin{equation}
    0,\quad
    \frac{\pi}{2},\quad
    \pi,\quad
    \frac{3\pi}{2},\quad
    2\pi.
\end{equation}

Do not give them the sine and cosine values immediately.

Instead, repeatedly ask:

\begin{quote}
    ``What are the coordinates of the point?''
\end{quote}

and then:

\begin{quote}
    ``If the coordinates are $(\cos\alpha,\sin\alpha)$,
    what are $\cos\alpha$ and $\sin\alpha$?''
\end{quote}

The completed table is

\begin{center}
\begin{tabular}{c|c|c}
    $\boldsymbol{\alpha}$ &
    $\boldsymbol{\cos(\alpha)}$ &
    $\boldsymbol{\sin(\alpha)}$ \\ \hline
    $0$              & $1$  & $0$  \\
    $\frac{\pi}{2}$  & $0$  & $1$  \\
    $\pi$             & $-1$ & $0$  \\
    $\frac{3\pi}{2}$ & $0$  & $-1$ \\
    $2\pi$            & $1$  & $0$
\end{tabular}
\end{center}


\subsubsection*{Preparing the Next Discovery}

Finally, ask the students to imagine that the point moves continuously
around the unit circle.

\begin{quote}
    ``While the point moves around the circle, what happens to its
    $y$-coordinate?''
\end{quote}

Collect a few observations:

\begin{equation}
    0
    \rightarrow
    1
    \rightarrow
    0
    \rightarrow
    -1
    \rightarrow
    0.
\end{equation}

Discuss the following questions in groups: 

\begin{quote}
    ``What would happen if we plotted this $y$-coordinate against the
    angle $\alpha$?''
\end{quote}





\section{Trigonometric function}

So far, we have considered individual angles and determined their sine
and cosine using the unit circle.
Now we want to change our perspective.

\begin{quote}
    ``What happens to sine and cosine if we do not look at just one angle,
    but continuously change the angle?''
\end{quote}

\subsection{Turning values into graphs}

\begin{align*}
    \alpha=0
        &\quad\Rightarrow\quad \sin(\alpha)=0,\\
    \alpha=\frac{\pi}{2}
        &\quad\Rightarrow\quad \sin(\alpha)=1,\\
    \alpha=\pi
        &\quad\Rightarrow\quad \sin(\alpha)=0,\\
    \alpha=\frac{3\pi}{2}
        &\quad\Rightarrow\quad \sin(\alpha)=-1,\\
    \alpha=2\pi
        &\quad\Rightarrow\quad \sin(\alpha)=0.
\end{align*}

Ask:

\begin{quote}
    ``What happens between these points?''
\end{quote}


The point on the circle moves continuously, so its vertical coordinate
also changes continuously.

Collect ideas, finally shpw simulation 

Simulation
https://www.geogebra.org/m/jfrrtwrw?utm_source=chatgpt.com

https://animg.app/en/animation/unit-circle-generates-sine-wave-hRwsY8Hw?utm_source=chatgpt.com


\begin{center}
\fbox{\parbox{0.82\textwidth}{
\centering
The sine function describes the vertical coordinate of a point
moving around the unit circle.
}}
\end{center}


\subsection*{ What Do We Notice About the Sine Function?}

Look at the resulting graph together.


\textbf{Discussion:}

\begin{itemize}
    \item What is the largest value the sine can have?
    \item What is the smallest value?
    \item At which angles is $\sin(\alpha)=0$?
    \item Can we continue rotating the point after one complete turn?
    \item Can we also rotate the point in the opposite direction?
    \item What happens after one complete turn around the circle?
\end{itemize}

\paragraph{Domain}

Ask:

\begin{quote}
    ``Which angles can we put into the sine function?''
\end{quote}

Using the unit circle, we can rotate the point by any angle.
We can continue for more than one full rotation, and we can also rotate
in the opposite direction, corresponding to negative angles.

Therefore, the sine function is defined for all real numbers:

\begin{equation}
    \boxed{\text{Domain: } \alpha \in \mathbb{R}}
\end{equation}

\paragraph{Range}

Now ask:

\begin{quote}
    ``Which values can come out of the sine function?''
\end{quote}

Since $\sin(\alpha)$ is the vertical coordinate of a point on the unit
circle, it can never be larger than $1$ or smaller than $-1$.

Therefore,

\begin{equation}
    \boxed{\text{Range: } -1 \leq \sin(\alpha) \leq 1}
\end{equation}

or, using interval notation,

\begin{equation}
    \boxed{\text{Range: } [-1,1]}.
\end{equation}

\paragraph{Periodicity}

Finally, consider what happens after one complete rotation.

After an angle of $2\pi$, the point has returned to exactly the same
position on the unit circle. Therefore, all sine values start repeating:

\begin{equation}
    \boxed{\sin(\alpha+2\pi)=\sin(\alpha)}.
\end{equation}

A function whose values repeat after a certain interval is called a
\textbf{periodic function}.

The sine function has the period

\begin{equation}
    \boxed{T=2\pi}.
\end{equation}

\paragraph{Summary}

From the unit circle and the graph, we have discovered fot the function $f(\alpha)= sin(\alpha)$:

\begin{itemize}
    \item \textbf{Domain:} $\mathbb{R}$
    \item \textbf{Range:} $[-1,1]$
    \item \textbf{Maximum:} $1$
    \item \textbf{Minimum:} $-1$
    \item \textbf{Period:} $2\pi$
\end{itemize}
\subsection*{ What About Cosine?}

Now return to the unit circle.

Ask:

\begin{quote}
    ``For the sine function, we followed the vertical coordinate.
    What would we have to follow to obtain the cosine function?''
\end{quote}

Since

\begin{equation}
    \cos(\alpha)=x_P,
\end{equation}

the cosine describes the \textbf{horizontal position} of the point on
the unit circle.

Again consider the important angles:

\begin{center}
\begin{tabular}{c|ccccc}
$\alpha$
& $0$
& $\frac{\pi}{2}$
& $\pi$
& $\frac{3\pi}{2}$
& $2\pi$ \\ \hline
$\cos(\alpha)$
& $1$
& $0$
& $-1$
& $0$
& $1$
\end{tabular}
\end{center}

Ask the students to predict the shape of the cosine graph before
showing or drawing it.

Then draw the sine and cosine functions in the same coordinate system.

Ask:

\begin{itemize}
    \item What do sine and cosine have in common?
    \item What is different?
    \item Do they have the same maximum and minimum values?
    \item Do they have the same period?
    \item Could we move one graph horizontally so that it lies exactly
    on top of the other?
\end{itemize}

The students should notice that both functions have the same shape but
start at different points.

The cosine function starts at its maximum,

\begin{equation}
    \cos(0)=1,
\end{equation}


while the sine function starts at a zero crossing,

\begin{equation}
    \sin(0)=0.
\end{equation}

The two graphs are shifted relative to each other by a quarter of a
complete period:

\begin{equation}
    \boxed{
    \cos(\alpha)
    =
    \sin\left(\alpha+\frac{\pi}{2}\right)
    }.
\end{equation}

This horizontal displacement is called a \textbf{phase shift}.

\subsection*{Transforming Sine and Cosine Functions}

So far, we have considered the basic sine and cosine functions

\begin{equation}
    f(x)=\sin(x)
    \qquad \text{and} \qquad
    f(x)=\cos(x).
\end{equation}

Both oscillate between $-1$ and $1$ and have a period of $2\pi$.

But real oscillations can look very different.

\textbf{Discussion:}

\begin{itemize}
    \item Do all oscillations have the same maximum displacement?
    \item Do all oscillations take the same amount of time for one cycle?
    \item Do all oscillations start at the same point?
    \item How could we modify our sine and cosine functions to describe
    these different situations?
\end{itemize}

We will investigate this by changing the functions step by step.
-
\subsubsection*{1. Changing the Amplitude}

Start with

\begin{equation}
    f(x)=\sin(x).
\end{equation}

Compare it with

\begin{equation}
    f(x)=2\sin(x)
\end{equation}

and

\begin{equation}
    f(x)=\frac{1}{2}\sin(x).
\end{equation}

Plot the three functions in the same coordinate system.

\textbf{Discussion:}

\begin{itemize}
    \item What changes?
    \item What stays the same?
    \item What are the maximum and minimum values of each function?
    \item Does the period change?
\end{itemize}

Let the students complete \textbf{Exercise 13}.

The students should discover that the number in front of the sine
changes how far the graph moves away from its middle line.

In general,

\begin{equation}
    f(x)=a\sin(x).
\end{equation}

The amplitude is

\begin{equation}
    \boxed{A=|a|}.
\end{equation}

For example,

\begin{equation*}
    3\sin(x)
\end{equation*}

oscillates between $-3$ and $3$.

If $a<0$, the graph is additionally reflected across the $x$-axis.


\subsubsection*{2. Moving the Whole Oscillation Up and Down}

Ask: How can you modify the function such that the whole graph is shifted up/down

\begin{equation}
    f(x)=\sin(x),
\end{equation}

\begin{equation}
    f(x)=\sin(x)+2,
\end{equation}

and

\begin{equation}
    f(x)=\sin(x)-1.
\end{equation}

Ask:

\begin{itemize}
    \item Does the shape of the graph change?
    \item Does the amplitude change?
    \item Does the period change?
    \item What happens to the middle line of the oscillation?
\end{itemize}

The students should discover that adding a number moves the entire
graph vertically.

In general,

\begin{equation}
    f(x)=\sin(x)+d.
\end{equation}

The middle line is now

\begin{equation}
    \boxed{y=d}.
\end{equation}

Therefore,

\begin{itemize}
    \item $d>0$: the graph moves upwards,
    \item $d<0$: the graph moves downwards.
\end{itemize}


% --------------------------------------------------
\subsubsection*{3. Changing the Period}

Return to the basic sine function and compare

\begin{equation}
    f(x)=\sin(x)
\end{equation}

with

\begin{equation}
    f(x)=\sin(2x).
\end{equation}

Let the students first look at the graphs without giving them the rule.

\textbf{Discussion:}

\begin{itemize}
    \item How many complete oscillations does $\sin(x)$ make between
    $0$ and $2\pi$?
    \item How many complete oscillations does $\sin(2x)$ make?
    \item Which function repeats faster?
    \item What is the period of $\sin(2x)$?
    \item Does the amplitude change?
\end{itemize}

Let the students complete \textbf{Exercise 14}.

Then consider another example:

\begin{equation}
    f(x)=\sin\left(\frac{1}{2}x\right).
\end{equation}

Ask the students to predict what will happen before plotting the function.

For the general function

\begin{equation}
    f(x)=\sin(bx),
\end{equation}

the period is

\begin{equation}
    \boxed{T=\frac{2\pi}{|b|}}.
\end{equation}

Thus, a larger value of $|b|$ means a shorter period and therefore
more oscillations in the same interval.

\begin{center}
\fbox{\parbox{0.82\textwidth}{
\centering
The factor outside the sine changes the \textbf{amplitude}.

The factor multiplying $x$ inside the sine changes the
\textbf{period}.
}}
\end{center}


% --------------------------------------------------
\subsubsection*{4. Shifting the Oscillation Horizontally}

Now compare

\begin{equation}
    f(x)=\sin(x)
\end{equation}

and

\begin{equation}
    f(x)=\sin\left(x-\frac{\pi}{2}\right).
\end{equation}

Ask:

\begin{itemize}
    \item Has the shape of the graph changed?
    \item Has the amplitude changed?
    \item Has the period changed?
    \item Where is the first graph at $x=0$?
    \item Where is the second graph at $x=0$?
    \item In which direction has the graph moved?
\end{itemize}

The students should discover that the entire oscillation has been
shifted horizontally.

For

\begin{equation}
    f(x)=\sin(x+c),
\end{equation}

the graph is shifted by

\begin{equation}
    -c.
\end{equation}

Therefore,

\begin{itemize}
    \item $\sin(x-c)$ is shifted to the \textbf{right} by $c$,
    \item $\sin(x+c)$ is shifted to the \textbf{left} by $c$.
\end{itemize}

This horizontal shift is called the \textbf{phase shift}.


% --------------------------------------------------
\subsubsection*{5. Putting Everything Together}

We can now combine all four transformations.

A general transformed sine function can be written as

\begin{equation}
    \boxed{
    f(x)=a\sin(bx+c)+d
    }.
\end{equation}

Similarly, for cosine,

\begin{equation}
    \boxed{
    f(x)=a\cos(bx+c)+d
    }.
\end{equation}

The parameters have the following effects:

\begin{center}
\begin{tabular}{c|c|c}
\textbf{Parameter} & \textbf{What does it change?} & \textbf{Value} \\ \hline
$a$ & Amplitude & $A=|a|$ \\[2mm]
$b$ & Period & $T=\dfrac{2\pi}{|b|}$ \\[3mm]
$c$ & Horizontal position / phase & shift $=-\dfrac{c}{b}$ \\[3mm]
$d$ & Vertical position & middle line $y=d$
\end{tabular}
\end{center}

\textbf{Important:} The same rules apply to both sine and cosine.


% --------------------------------------------------
\subsubsection*{6. Can We Read a Function?}

Before calculating anything, show the students an example:

\begin{equation}
    f(x)=2\sin\left(3x-\frac{\pi}{2}\right)+1.
\end{equation}

Ask them to identify the four parameters:

\begin{equation}
    a=2,
    \qquad
    b=3,
    \qquad
    c=-\frac{\pi}{2},
    \qquad
    d=1.
\end{equation}

Then determine the properties step by step:

\begin{align}
    \text{Amplitude:}\qquad
    A &= 2,\\
    \text{Period:}\qquad
    T &= \frac{2\pi}{3},\\
    \text{Phase shift:}\qquad
    -\frac{c}{b}
      &= \frac{\pi}{6},\\
    \text{Middle line:}\qquad
    y &= 1.
\end{align}

Therefore, the function describes an oscillation with amplitude $2$
around the middle line $y=1$, with a period of $2\pi/3$, shifted by
$\pi/6$ to the right.


% --------------------------------------------------
\subsubsection*{7. From Mathematics to Physics}

Finally, return to the question from the beginning:

\begin{quote}
    ``What physical systems or phenomena can be described by sine and cosine functions?''
\end{quote}


\textbf{Possible answers:}
\begin{itemize}
    \item a swinging pendulum
    \item a mass attached to a spring
    \item a vibrating string
    \item sound waves
    \item water waves (approximately)
    \item alternating current
    \item electromagnetic waves, such as light
\end{itemize}

\end{document}
